% Einheit 2 von 4 – Scheitelpunktform (bank/quadratische-funktionen/e2.jsonl, zusammenbau v0.1)
%% TODO Zweigzeile Teil 1: was hier gelernt wird (Satz in Schülersprache aus dem Einheitstitel) – Einheit 2
\einheitenkopf[e2]{Einheit 2 von 4 · Scheitelpunktform}
\zweigzeile{ab Kl. 9, je nach Buch bis Kl. 10 · P10 oft}
\setcounter{aufgabe}{12}

%% TODO Titel: Ich-kann-Satz fehlt in der Bank (Platzhalter: Kettenname) – Nr. 13
%% TODO Anweisung: ein Satz über der ersten Teilaufgabe fehlt in der Bank – Nr. 13
\begin{aufgabe}{Scheitelpunktform}
%% TODO \rechenplatz{n}: Schrittzahl des Grundfalls fehlt in der Bank (nur bei fester Schreibform, 2.3 b)
\begin{teile}
\teil $f(x) = (x - 4)^2 + 2$ – kreise die Zahl in der Klammer ein. Bei welchem $x$ liegt der Scheitel? x = \leerfeld
\teil $f(x) = (x - 2)^2 + 5$ – Scheitel? S( \leerfeld | \leerfeld )
\teil $f(x) = (x + 4)^2 + 2$ – Scheitel? S( \leerfeld | \leerfeld )
\teil Lies den Scheitel der Parabel ab.

\begin{ksys}[xmin=-5,xmax=1,ymin=-1,ymax=8,ablesen]
\parabel{1}{-2}{3}{f}
\end{ksys}
S( \leerfeld | \leerfeld )
\teil $f(x) = (x - 2)^2 + 1$ – skizziere die Parabel.

\begin{ksys}[xmin=-2,xmax=6,ymin=-1,ymax=10]
\end{ksys}
\teil Nach oben geöffnete Normalparabel mit Scheitel $S(5|2)$ – Gleichung? f(x) = \leerfeld
\teil Gleichung der Parabel?

\begin{ksys}[xmin=-1,xmax=5,ymin=-4,ymax=2,ablesen]
\parabel{1}{2}{-3}{f}
\end{ksys}
f(x) = \leerfeld
\teil $f(x) = (x - 1)^2 + 2$ wird um $4$ nach unten verschoben – Gleichung der neuen Parabel $g$? g(x) = \leerfeld
\teil $f(x) = (x - 5)^2 + 1$ wird an der $x$-Achse gespiegelt – Gleichung des Spiegelbilds $g$? g(x) = \leerfeld
\teil $f(x) = (x - 4)^2 + 2$ wird an der $y$-Achse gespiegelt – Gleichung des Spiegelbilds $g$? g(x) = \leerfeld
\teil Gib eine nach oben geöffnete Normalparabel an, deren Scheitel auf der $x$-Achse rechts vom Ursprung liegt. f(x) = \leerfeld
\teil $f(x) = (x - 1)^2 + 2$ und $g(x) = (x - 1)^2 - 4$ – wie liegen die Parabeln zueinander? Haben sie gemeinsame Punkte? Begründe ohne Rechnung.
\teil $f(x) = 2(x - 4)^2 + 1$ – Scheitel, Öffnung, schmaler oder breiter als die Normalparabel? Punkt eine Einheit rechts vom Scheitel? S( \leerfeld | \leerfeld ), Punkt ( \leerfeld | \leerfeld )
\teil $f(x) = (x - 4)^2 - 5$ wird erst um $5$ nach oben verschoben und dann an der $x$-Achse gespiegelt. Gleichung der neuen Parabel $g$? \hfill (P10 2017 OS) \\ g(x) = \leerfeld
\teil Die Parabel $f(x) = (x + 4)^2 - 2$ wird an der $y$-Achse gespiegelt. Skizziere das Spiegelbild $g$ und gib seine Gleichung an. \hfill (P10 2020 OS) \\

\begin{ksys}[xmin=-7,xmax=7,ymin=-3,ymax=7]
\parabel{1}{-4}{-2}{f}
\end{ksys}
g(x) = \leerfeld
\teil Skizziere eine nach unten geöffnete Normalparabel $f$ mit Scheitel auf der $y$-Achse, die keine Nullstelle hat. Gib ihre Gleichung und die Gleichung ihres Spiegelbilds $g$ an der $x$-Achse an. \hfill (P10 2018 OS) \\

\begin{ksys}[xmin=-4,xmax=4,ymin=-10,ymax=10,ystep=2]
\end{ksys}
f(x) = \leerfeld , g(x) = \leerfeld
\end{teile}
\end{aufgabe}

%% TODO Titel: Ich-kann-Satz fehlt in der Bank (Platzhalter: Kettenname) – Nr. 14
%% TODO Anweisung: ein Satz über der ersten Teilaufgabe fehlt in der Bank – Nr. 14
\begin{aufgabe}{Scheitelpunktform – Fehler finden, Begründen, Darstellungswechsel, Anwendung}
\begin{teile}
\teil Ben liest bei $f(x) = (x + 6)^2 - 2$ den Scheitel $S(6|{-2})$ ab. Finde den Fehler und korrigiere.
\teil Warum hat $f(x) = (x - 4)^2 + 1$ bei $x = 4$ ihren kleinsten Wert?
\teil Wahr oder falsch? \\ Die Gleichung ist $f(x) = (x + 1)^2 - 4$. \kreuz{wahr} \kreuz{falsch} \\ Der Scheitel ist $S(1|{-4})$. \kreuz{wahr} \kreuz{falsch} \\ Die Parabel ist nach unten geöffnet. \kreuz{wahr} \kreuz{falsch}

\begin{ksys}[xmin=-2,xmax=4,ymin=-5,ymax=1,ablesen]
\parabel{1}{1}{-4}{f}
\end{ksys}
\teil Beim Kugelstoßen gilt für die Flugbahn $h(x) = -0{,}05(x - 6)^2 + 3{,}8$ ($x$: Entfernung in m, $h$: Höhe in m). Wie hoch fliegt die Kugel höchstens, und in welcher Entfernung?

\begin{ksys}[xmin=0,xmax=16,ymin=0,ymax=5,xstep=2,xlabel=$x$ in m,ylabel=$h$ in m]
\funktionab{-0.05*(\x-6)^2+3.8}{h}{0}{14.7}
\end{ksys}
\leerfeld[m] hoch, \leerfeld[m] entfernt
\end{teile}
\end{aufgabe}
